% FORMALIZACION_NUEVA de la incidencia firmada de la carta Paley preexistente.
% La construcción comienza en exc:aw; no identifica por decreto sus banderas
% con el cociente de todas las historias TPK.
\paragraph{Construction du porteur à partir de l'incidence à six positions.}
\label{ii:ico:construccion-firmada}
La matrice \(A_W\) de \eqref{exc:aw}, qui intervient préalablement dans le
transport angulaire de la section~\ref{sec:ii-enlace-angular-witt}, fournit
un antécédent explicite de la réalisation icosaédrique. Soit
\(C=\operatorname{bal}(A_W)\), où
\(\operatorname{bal}(0)=0\), \(\operatorname{bal}(1)=1\) et
\(\operatorname{bal}(2)=-1\). L’ordre de ses coordonnées est
\(\mathcal U=(\infty,0,1,2,3,4)\). La coordinatisation entière satisfait
\(C=C^{\mathsf T}\), \(C^2=5I_6\) et \(\operatorname{Tr}C=0\).
Nous définissons son revêtement signé par
\begin{equation}
 \mathcal D_C=\mathcal U\times\{+1,-1\},\qquad
 (u,\sigma)\sim_C(v,\tau)
 \ \Longleftrightarrow\ u\ne v\ \text{ et }\ \sigma\tau=C_{uv}.
 \label{ii:ico:incidencia-firmada}
\end{equation}
L'orientation dédouble chaque coordonnée ; la matrice décide des incidences
entre les deux fibres correspondantes. L'involution de ce revêtement
est \(\jmath_C(u,\sigma)=(u,-\sigma)\).

\Needspace{12\baselineskip}
\begin{samepage}
\begin{proposition}[Réalisation isométrique de l'incidence signée]
\label{ii:ico:realizacion-firmada}
L'opérateur et les vecteurs
\begin{equation}
 E_C=\frac12\left(I_6+\frac C{\sqrt5}\right),\qquad
 \mathcal H_C=\operatorname{im}E_C,\qquad
 v_{u,\sigma}=\sigma\sqrt2 E_Ce_u
 \label{ii:ico:proyector-seis}
\end{equation}
réalisent \(\mathcal D_C\) comme douze vecteurs unitaires distincts dans un
espace euclidien de dimension trois. Il existe une isométrie explicite
\(T_C:\mathcal H_C\to\mathbb R^3\) pour laquelle
\(\iota_C(u,\sigma)=T_Cv_{u,\sigma}\) est une bijection sur
\(\Omega_{12}^{\rm ico}\), conserve l'incidence et satisfait
\(\iota_C\jmath_C=-\iota_C\).
\end{proposition}
\end{samepage}
\begin{proof}
De \(C^2=5I_6\), on obtient \(E_C^2=E_C=E_C^*\) ; sa trace est trois.
Les identités métriques sont
\[
 \langle v_{u,\sigma},v_{v,\tau}\rangle
 =\begin{cases}\sigma\tau,&u=v,\\
 \sigma\tau C_{uv}/\sqrt5,&u\ne v.
 \end{cases}
\]
Ainsi, les douze vecteurs sont distincts et leur incidence de produit scalaire
\(1/\sqrt5\) est exactement~\eqref{ii:ico:incidencia-firmada}.
Pour expliciter l'isométrie, prenons pour colonnes de \(B\), dans l'ordre
\(\mathcal U\),
\[
 B=\frac1\nu
 \begin{pmatrix}
 0&-1&-\varphi_g&0&\varphi_g&1\\
 -1&-\varphi_g&0&1&0&-\varphi_g\\
 -\varphi_g&0&-1&-\varphi_g&-1&0
 \end{pmatrix}.
\]
La substitution de \(\varphi_g^2=\varphi_g+1\) donne
\(B^*B=I_6+C/\sqrt5=2E_C\). Ainsi
\(T_C=B/\sqrt2\big|_{\mathcal H_C}\) est une isométrie, et
\(T_Cv_{u,\sigma}=\sigma Be_u\). Les six colonnes de \(B\) et leurs
opposées sont précisément les douze sommets de
\eqref{ii:pent:eq:gravity-icosahedron-vertices}. La dernière identité
démontre également la compatibilité des involutions.
\end{proof}

La réalisation dodécaphasique précédente construit les douze vecteurs
directement à partir du projecteur \(P_3^{\mathrm{ico}}\) et des positions du calendrier.
Le revêtement signé fournit une autre expression en coordonnées de ce
porteur. Pour le lecteur \(\pi_{\rm ico}\) du théorème de descente,
la relation entre les deux expressions s'écrit
\begin{equation}
 X_{\rm TPK}^{\rm exc}\xrightarrow{\ r_\pm\ }
 \mathcal D_C\xrightarrow[\simeq]{\ \iota_C\ }
 \Omega_{12}^{\rm ico},\qquad
 \pi_{\rm ico}(x)=\sigma(x)Be_{u(x)},\quad
 r_\pm(x)=(u(x),\sigma(x)).
 \label{ii:ico:factorizacion-lector}
\end{equation}
Comme \(\iota_C\) est une bijection, cette factorisation détermine
\(r_\pm=\iota_C^{-1}\pi_{\rm ico}\). Par conséquent, \(r_\pm\) exprime
dans la coordinatisation signée le lecteur utilisé : son introduction
n’ajoute pas d’hypothèse indépendante pour construire le porteur
dodécaphasique. Le tableau de correspondances précédent donne cette expression
pour ses douze positions. Les propriétés (i)--(iv) conservent leur
portée propre lorsque l’on identifie la réalisation au quotient
du chemin exceptionnel complet. En particulier, le transport conjoint
du repère, du projecteur et de l’incidence démontre la covariance entre coordinatisations ;
la conservation d’une incidence fixe par une transformation active
est une identité distincte. L’orientation et la mémoire restent
dans l’état source pendant les deux lectures.
Une permutation signée qui transporte \(C\) transporte conjointement
\(E_C\), ses vecteurs et son incidence ; cette covariance permet de comparer
les coordinatisations en conservant l’orientation de chaque fibre.
